\documentclass[aspectratio=169,11pt]{beamer}
\usepackage{iftex}
\ifPDFTeX
  \usepackage[T1]{fontenc}
\else
  \usepackage{fontspec}
\fi
\usepackage{tikz}
\usetikzlibrary{positioning,arrows.meta}

\newif\ifusebrandtheme
\usebrandthemefalse
\ifdefined\UseUnofficialRedHatTheme
  \ifXeTeX
    \IfFileExists{beamerthemeredhat.sty}{
      \IfFileExists{img/logo.png}{
        \IfFileExists{img/logo_white.png}{
          \IfFontExistsTF{Red Hat Text}{
            \IfFontExistsTF{Red Hat Mono}{\usebrandthemetrue}{}
          }{}
        }{}
      }{}
    }{}
  \fi
  \ifusebrandtheme\else
    \PackageWarning{libvirt-sno-testing}{Unofficial theme prerequisites missing; using neutral theme}
  \fi
\fi

\ifusebrandtheme
  \usetheme{redhat}
\else
  \usetheme{Madrid}
  \definecolor{labblue}{RGB}{28,65,89}
  \usecolortheme[named=labblue]{structure}
  \setbeamertemplate{blocks}[rounded][shadow=false]
  \setbeamertemplate{title page}[default][rounded=true,shadow=false]
  \setbeamertemplate{footline}[frame number]
\fi
\setbeamertemplate{navigation symbols}{}
\setbeamersize{text margin left=9mm,text margin right=9mm}
\setbeamertemplate{itemize items}[circle]

\newcommand{\reposource}[1]{\par\vfill{\tiny Source: #1\par}}
\newcommand{\readmeref}[1]{\href{https://github.com/matejvasek/mycluster/blob/master/README.md}{README: #1}}
\newcommand{\scriptref}[1]{\href{https://github.com/matejvasek/mycluster/blob/master/deploy-sno.sh}{deploy-sno.sh: #1}}

\title{An OpenShift Test Lab on Libvirt}
\subtitle{Single-node provisioning, checks, and operational trade-offs}
\author{}
\date{}

\begin{document}

\begin{frame}
  \titlepage
  \begin{center}
    \small Agent-based SNO installation with \texttt{matejvasek/mycluster}\\[1mm]
    The repository provisions the lab; workload tests live elsewhere.
  \end{center}
\end{frame}

\begin{frame}{Topology: one VM, routed dual-stack networking}
  \centering
  \begin{tikzpicture}[
    box/.style={draw=structure.fg,rounded corners,align=center,
      minimum height=15mm,text width=33mm,font=\small},
    >={Latex}]
    \node[box] (lan) {LAN / upstream router\\routes VM subnets};
    \node[box,right=8mm of lan] (host) {KVM host\\libvirt bridge + dnsmasq};
    \node[box,right=8mm of host] (vm) {SNO VM\\static IPv4 + IPv6};
    \draw[<->,thick] (lan) -- (host);
    \draw[<->,thick] (host) -- (vm);
  \end{tikzpicture}
  \begin{itemize}
    \item Routed network, not NAT: upstream routes point to the host.
    \item Agent-based installer; one control-plane node, zero workers;
      \texttt{platform: none} and OVN-Kubernetes.
    \item Two qcow2 disks: OS plus data for optional LVM Storage; UEFI boot.
    \item \textbf{Defaults:} 12 vCPUs, 32 GiB RAM, 120 GiB OS, 80 GiB data.
  \end{itemize}
  \reposource{\readmeref{Network routing, Configuration}; \scriptref{generate\_configs, create\_vm}}
\end{frame}

\begin{frame}{Inputs: adapt the defaults to the host and LAN}
  \begin{itemize}
    \item Tools: \texttt{openshift-install}, \texttt{oc}, \texttt{virsh},
      \texttt{virt-install}, \texttt{nmstatectl}; host administrative access
      for libvirt/network operations.
    \item Local pull secret and an SSH \emph{public} key via \texttt{SSH\_KEY}.
    \item Configure \texttt{CLUSTER\_NAME}, \texttt{BASE\_DOMAIN},
      \texttt{IPV4\_SUBNET}, \texttt{IPV6\_PREFIX}, and VM sizing.
      These settings accept environment overrides; optional \texttt{./.env}
      is sourced first and can itself set values.
    \item Optional TLS chain/key: API and application-route wildcard SANs;
      without the files, custom TLS configuration is skipped.
    \item README cautions against allocating much more than half the host's
      RAM to the VM: leave headroom for host services and page cache.
  \end{itemize}
  \begin{block}{Keep local inputs local}
    \small \texttt{.gitignore} excludes the pull secret, TLS files,
    \texttt{.env}, install directory, and ISO cache. Ignoring is not encryption.
  \end{block}
  \reposource{\readmeref{Prerequisites, Configuration, Caveats}; \scriptref{configuration, preflight};
    \href{https://github.com/matejvasek/mycluster/blob/master/.gitignore}{.gitignore}}
\end{frame}

\begin{frame}[fragile]{Deployment: the documented sequence}
\begin{verbatim}
./deploy-sno.sh
\end{verbatim}
  \begin{enumerate}
    \item \texttt{preflight}: check tools, pull secret, SSH key; inspect TLS SANs.
    \item \texttt{create\_network}: routed libvirt network, DNS, host resolver.
    \item \texttt{generate\_configs}: install-config and agent-config YAML.
    \item \texttt{generate\_iso}: agent image; cache keyed by configs + installer version.
    \item \texttt{create\_vm}: copy ISO into libvirt images; boot VM with two disks.
    \item \texttt{wait\_for\_install}: bootstrap, then install completion.
    \item \texttt{apply\_certs}: optional custom ingress and API certificates.
  \end{enumerate}
  \small Individual functions are callable as \texttt{./deploy-sno.sh <function>}.
  \alert{Re-running the full deployment replaces the VM and both disks.}
  \reposource{\readmeref{Usage}; \scriptref{main, generate\_iso, create\_vm}}
\end{frame}

\begin{frame}[fragile]{After install: smoke checks, not workload tests}
\begin{verbatim}
export KUBECONFIG=./ocp-sno-install/auth/kubeconfig
oc get nodes
oc get clusterversion
oc get clusteroperators
\end{verbatim}
  \begin{itemize}
    \item Suggested interpretation: inspect node readiness, cluster version,
      and operator conditions before using the lab.
    \item Console: \texttt{https://console-openshift-console.}\\
      \texttt{apps.<cluster>.<domain>}
    \item These are the documented post-install checks, not a test suite.
    \item Workload tests are outside this repository; choose and run those
      separately against the provisioned cluster.
  \end{itemize}
  \reposource{\readmeref{After installation}; \scriptref{wait\_for\_install, main}}
\end{frame}

\begin{frame}[fragile]{Optional: persistent storage and an internal registry}
  \begin{columns}[T,onlytextwidth]
    \begin{column}{0.48\textwidth}
      \begin{block}{1. Install LVM Storage}
\begin{verbatim}
./deploy-sno.sh install_lvms
\end{verbatim}
        \small Uses the second data disk for dynamic PV provisioning.
        Selects the operator channel from the OCP version and makes
        \texttt{lvms-vg1} the default StorageClass.
      \end{block}
    \end{column}
    \begin{column}{0.48\textwidth}
      \begin{block}{2. Enable the registry}
\begin{verbatim}
./deploy-sno.sh enable_registry
\end{verbatim}
        \small Disabled by default on \texttt{platform: none}.
        Creates a 20 GiB RWO PVC on \texttt{lvms-vg1}; one replica
        with a \texttt{Recreate} rollout strategy.
      \end{block}
    \end{column}
  \end{columns}
  \medskip
  \small These are explicit post-install steps, not part of \texttt{main}.
  The registry configuration depends on the LVMS StorageClass.
  \reposource{\readmeref{Post-install optional steps, Notes}; \scriptref{install\_lvms, enable\_registry}}
\end{frame}

\begin{frame}{Networking and DNS: the easy places to get stuck}
  \begin{itemize}
    \item Upstream routes: IPv4 VM subnet via the host's LAN IPv4 address;
      IPv6 /64 via the host's link-local IPv6 address.
    \item dnsmasq serves \texttt{api}, \texttt{api-int}, and \texttt{*.apps};
      \texttt{setup\_dns} routes host base-domain queries to the bridge DNS.
    \item Other clients need cluster records in their DNS, pointing at node IPs.
    \item \alert{No DNS wildcard at \texttt{*.<cluster>.<domain>}.}
      \texttt{validateNoWildcardDNS.<cluster>.<domain>} must return NXDOMAIN.
      Application wildcard DNS and wildcard TLS SANs are distinct concerns.
    \item The README also warns that a cluster-domain wildcard breaks pod
      short-name resolution through search domains and \texttt{ndots:5}.
  \end{itemize}
  \reposource{\readmeref{Network routing, DNS, TLS certificates}; \scriptref{create\_network, setup\_dns}}
\end{frame}

\begin{frame}[fragile]{Recovery: distinguish slow installation from a broken lab}
  \begin{itemize}
    \item Install completion may time out around 70\%; README recommends
      retrying the wait step, rather than rebuilding immediately:
  \end{itemize}
\begin{verbatim}
./deploy-sno.sh wait_for_install
\end{verbatim}
  \begin{itemize}
    \item During installation the script restarts a VM that shuts off
      instead of rebooting.
    \item Host resolver settings are transient; after a host reboot:
  \end{itemize}
\begin{verbatim}
./deploy-sno.sh setup_dns
\end{verbatim}
  \begin{itemize}
    \item After more than 24 hours powered off, kubelet client certificates
      may have expired. README describes pending CSR approval, possibly
      in two rounds. \textbf{Suggestion:} inspect and approve only expected
      node CSRs, rather than approving all requests blindly.
  \end{itemize}
  \reposource{\readmeref{Caveats, Notes}; \scriptref{wait\_for\_install, setup\_dns}}
\end{frame}

\begin{frame}[fragile]{Lifecycle: preserve state or deliberately rebuild}
  \textbf{Keep the VM and disks} (documented default VM name):
\begin{verbatim}
sudo virsh shutdown ocp-sno
sudo virsh start ocp-sno
\end{verbatim}
  \begin{itemize}
    \item \textbf{Suggestion:} confirm the guest is shut off before stopping
      the host; re-run smoke checks after startup.
    \item \textbf{Destructive cleanup:} README lists VM destroy/undefine
      (including NVRAM), removal of both qcow2 disks and the copied ISO,
      then network destroy/undefine. Adapt names and paths if overridden.
    \item A fresh full deployment recreates the VM and disks; the ISO cache
      can avoid image generation when configs and installer match.
    \item \textbf{Suggestion:} back up needed workloads before cleanup or
      re-running \texttt{create\_vm}; a cached ISO is not a workload backup.
  \end{itemize}
  \reposource{\readmeref{Cluster lifecycle, ISO caching}; \scriptref{create\_vm, generate\_iso}}
\end{frame}

\begin{frame}{Takeaway: a provisioned lab is the starting point}
  \begin{block}{What this repository supplies}
    A script to provision dual-stack SNO on libvirt, documented smoke
    checks, and optional storage, registry, TLS, and identity helpers.
  \end{block}
  \begin{block}{What to plan separately}
    Workload tests, their success criteria, and any workload backup or
    restore procedure. No workload test suite is included here.
  \end{block}
  \begin{itemize}
    \item Suggested checklist: size the host, verify routes/DNS, deploy,
      inspect smoke checks, then run your separately chosen tests.
    \item Know whether the next action preserves disks or recreates them.
  \end{itemize}
  \reposource{\href{https://github.com/matejvasek/mycluster}{Public repository};
    \readmeref{Usage, Cluster lifecycle}; \scriptref{main, optional functions}}
\end{frame}

\end{document}
